\documentclass[11pt]{article}
\usepackage[T1]{fontenc}
\usepackage[margin=0.95in]{geometry}
\usepackage{amsmath,amssymb,amsthm,booktabs,array,hyperref,needspace}
\hypersetup{colorlinks=true,linkcolor=blue,citecolor=blue,urlcolor=blue}
\emergencystretch=1.5em
\newtheorem{theorem}{Theorem}
\newtheorem{lemma}[theorem]{Lemma}
\newtheorem{proposition}[theorem]{Proposition}
\newtheorem{corollary}[theorem]{Corollary}
\newcommand{\Tr}{\operatorname{Tr}}
\newcommand{\co}{\operatorname{co}}
\newcommand{\Smin}{S_{\min}}
\title{Classical capacities and complementary simulation\\for qubit Davies channels}
\author{Lluis Eriksson}
\date{17 October 2026 --- Working paper}
\begin{document}\maketitle
\begin{abstract}
We extend the thermal block entropy estimate for generalized amplitude damping to the full qubit Davies region, including additional dephasing. The proof combines the triangular inequality of Tang, Zhu, Bai and Wang with a controlled-purification consequence of strong subadditivity: a third vertex completes the fixed-entropy decomposition used for the thermal boundary. The resulting population-resolved bound holds with every finite-dimensional partner. It implies strong superadditivity of the convex roof of output entropy, and hence arbitrary-partner additivity of Holevo information and classical capacity for both the Davies channel and its complementary channel. We give scalar convex-envelope formulas for both capacities. Every bipartite state supported on the two-dimensional Stinespring range has strongly additive entanglement of formation, so its state entanglement cost requires no regularization. A scalar moment program evaluates that formation functional for arbitrary input coherence. The state/channel correspondence, complement duality and existing axial convex-roof geometry are explicitly attributed to prior work. A normalized reverse channel stays within the Davies region, allowing us also to remove regularization from the parallel entanglement cost of the complementary channel. We obtain an elementary rate in a known flat-roof region and exhibit a nonunital example whose simulation cost strictly exceeds the cost of its Choi state. Finite-block and numerical bounds support the results. The triangular theorem remains an external dependency; the original channel's simulation cost, adaptive simulation and quantum capacity are not determined.
\end{abstract}

\section{Channel family, antecedents and result}
A qubit can relax toward a thermal population while its coherence decays faster than relaxation alone requires. Write the corresponding channel as
\begin{equation}\label{eq:channel}
 \Phi_{a,b,\lambda}
 \begin{pmatrix}1-p&z\\\bar z&p\end{pmatrix}
 =
 \begin{pmatrix}1-t(p)&\lambda z\\\lambda\bar z&t(p)\end{pmatrix},
 \qquad t(p)=a+kp,\quad k=1-a-b,
\end{equation}
where
\begin{equation}\label{eq:region}
 a,b\ge0,\qquad a+b\le1,\qquad 0\le\lambda^2\le k,\qquad\lambda\ge0.
\end{equation}
A fixed phase of the coherence multiplier can be removed by an output unitary. The map is completely positive: it is generalized amplitude damping followed by dephasing with multiplier $\lambda/\sqrt{k}$ when $k>0$. The limiting case $k=0$ is a constant diagonal-output channel.

The region \eqref{eq:region} is the closure of the usual dissipative qubit Davies family, allowing degenerate limits and exchange of the energy labels. Roga, Fannes and \.{Z}yczkowski characterized these maps and their single-use minimum output entropy \cite{roga}. At finite positive relaxation times their parameters satisfy
\[
 k=e^{-t/T_1},\qquad \lambda=e^{-t/T_2},\qquad T_2\le2T_1.
\]
Here $t$ in the exponents is time, whereas $t(p)$ in \eqref{eq:channel} is an output population. We use \eqref{eq:channel} as the mathematical definition throughout. General phase-covariant qubit channels allow the larger complete-positivity region $\lambda^2\le k+ab$; that entire region is not claimed here.

The proof builds on two recent results. Tang et al.\ established strong additivity for qubit channels admitting a pure output, using a triangular block entropy inequality \cite{tang}. OpenAI extended that inequality to the thermal equality case and obtained strong additivity for generalized amplitude damping \cite{openai}. In our parameters that boundary is $\lambda^2=k$. We add a third vertex to their fixed-entropy decomposition, proving the inequality for $\lambda^2<k$ as well. The new vertex follows directly from strong subadditivity, and is not a new entropy axiom. The scalar convex envelope is necessary because extra dephasing need not preserve convexity of the pure-output entropy.

Several overlaps are already understood: unital qubit channels \cite{king}, entanglement-breaking channels \cite{shor}, and the pure-output boundary \cite{tang}. Earlier work treated approximate Holevo additivity for damping combined with dephasing \cite{leditzky}; its ordinary damping examples now also fall within the pure-output theorem. Our contribution is the stated block extension and its additivity consequences across the thermal interior. Roga explicitly listed Davies capacity additivity as open in 2011 \cite[Section 6, p.~78]{rogathesis}; this historical statement is not a certificate of present priority. We do not claim a new description of Davies dynamics or a new single-qubit minimum-entropy optimization. The state formation and complementary-channel consequences below use the established Matsumoto--Shimono--Winter correspondence \cite{msw}, complement duality \cite{conjugate}, and the regularized formation characterization of state cost \cite{hht}. Hellmund and Uhlmann already analyzed the axial one-use convex roofs \cite{hu}; the new operational step is the stated strong additivity, not rediscovery of those roofs.


The simulation result concerns the complementary channel and the parallel
diamond-norm task of Berta et al.\ \cite{bcbw}. Wilde already determined the
cost of epolarizing channels, the complements of depolarizing channels, even
under sequential tests \cite[Section IV.D]{wilde}. More generally, the unital
qubit case is covered by the unital additivity and covariant channel/state
machinery \cite{king,msw,wilde}. We do not claim these cases as new. The
additional nonunital region is useful precisely because the optimizing
reference input need not be a Bell pair; Section~\ref{sec:simulation} gives
an exact example. The scalar roof and the parallel simulation theorem are
antecedents; closure of the explicit reverse parameters lets us apply the
Davies range theorem to this different bipartition.

All spaces below are finite dimensional. Let $S(\rho)=-\Tr\rho\ln\rho$ be entropy in nats and put
\[
 h(p)=-p\ln p-(1-p)\ln(1-p),\qquad
 g(u)=h\left(\frac{1+\sqrt{1-4u}}2\right).
\]
Zero terms are interpreted by continuity. Define
\begin{equation}\label{eq:ef}
 v(p)=t(p)(1-t(p))-\lambda^2p(1-p),\quad
 e(p)=g(v(p)),\quad f=\co e,
\end{equation}
where $\co e$ is the largest convex function below $e$ on $[0,1]$.

For a channel $\mathcal N$, $\chi(\mathcal N)$ denotes the maximum ensemble Holevo information in bits, $\Smin(\mathcal N)=\min_\rho S(\mathcal N(\rho))$, and
\[
 C(\mathcal N)=\sup_{n\ge1}\frac1n\chi(\mathcal N^{\otimes n})
\]
is its unassisted memoryless classical capacity, with arbitrary entangled block encodings and collective output measurements \cite{holevo,sw}. No preshared entanglement or feedback is assumed.

\begin{theorem}\label{thm:main}
For every channel $\Phi=\Phi_{a,b,\lambda}$ satisfying \eqref{eq:region},
\begin{equation}\label{eq:capacity}
 C(\Phi)=\chi(\Phi)=\frac1{\ln2}\max_{0\le p\le1}
       \{h(a+kp)-f(p)\}.
\end{equation}
An ensemble with at most four pure signals, consisting of at most two equally weighted opposite-phase pairs within their respective populations, attains the one-use maximum. For every finite-dimensional channel $\mathcal N$,
\begin{align}
 \chi(\Phi\otimes\mathcal N)&=\chi(\Phi)+\chi(\mathcal N),\label{eq:chiadd}\\
 \Smin(\Phi\otimes\mathcal N)&=\min_p e(p)+\Smin(\mathcal N),\label{eq:sadd}\\
 C(\Phi\otimes\mathcal N)&=\chi(\Phi)+C(\mathcal N).\label{eq:cadd}
\end{align}
In particular $\chi(\Phi^{\otimes n})=n\chi(\Phi)$ for every $n$.
\end{theorem}
Theorem~\ref{thm:main} states unrestricted Holevo additivity. The stronger fixed-marginal convex-roof statement is proved in Section~\ref{sec:formation}. The partner's classical capacity retains its regularization; we do not replace $C(\mathcal N)$ by $\chi(\mathcal N)$. Formula \eqref{eq:capacity} is a scalar convex-envelope optimization, not a claimed elementary closed form.

\section{The block entropy extension}
We first state precisely the external inequality used in the proof.

\begin{lemma}[Triangular entropy inequality, Tang et al.]\label{lem:triangular}
Let $A,B,C$ have the same finite rectangular shape, with squared Hilbert--Schmidt norms $x,y,z$ summing to one. Then
\[
 S\left(
 \begin{pmatrix}A&B\\C&0\end{pmatrix}
 \begin{pmatrix}A&B\\C&0\end{pmatrix}^{\!*}\right)
 \ge xS(AA^*/x)+yS(BB^*/y)+zS(CC^*/z)+g(yz).
\]
A zero-weight term is omitted.
\end{lemma}
This is the same-shape rectangular case of \cite[Theorem 4.2]{tang}, after permutation of block rows, also stated and proved in \cite[Theorem 2.1]{openai}. We invoke that result; it is not proved anew or attributed to this paper.

The additional vertex uses only data processing for conditional entropy. Its two-label case also follows from Amosov's phase-damping estimate \cite[Proposition 5]{amosov}, with partial trace as the partner channel. We give the short purification proof to fix the direction of the entropy inequality; the vertex estimate itself is not claimed as new.
\begin{lemma}[Controlled purifications]\label{lem:controlled}
Let $\rho=(\rho_{ij})$ be a density matrix and let $X_i$ be matrices of the same shape with $\Tr X_iX_i^*=1$. Then
\[
 S\left(\sum_{ij}\rho_{ij}|i\rangle\langle j|\otimes X_iX_j^*\right)
 \ge S(\rho)+\sum_i\rho_{ii}S(X_iX_i^*).
\]
\end{lemma}
\begin{proof}
Vectorize $X_i$ as a normalized vector $|X_i\rangle_{BE}$. The map
$V|i\rangle=|i\rangle_A|X_i\rangle_{BE}$ is an isometry. Thus
$\Omega=V\rho V^*$ has $S(ABE)_\Omega=S(\rho)$, and its $AB$ marginal is the matrix in the statement. Dephase $A$, obtaining $\Omega'$. Strong subadditivity implies that applying a channel to the conditioning system $AB$ cannot decrease $S(E|AB)$ \cite{ssa}. Consequently
\[
 S(AB)_\Omega-S(ABE)_\Omega
 \ge S(AB)_{\Omega'}-S(ABE)_{\Omega'}.
\]
The right side is $\sum_i\rho_{ii}S(X_iX_i^*)$: the classical entropy of $i$ cancels, and each conditional $BE$ vector is pure. Rearrangement proves the claim.
\end{proof}

\begin{theorem}[Thermal inequality with additional dephasing]\label{thm:block}
Let $X,Y\in\mathbb C^{d\times r}$ satisfy $\Tr XX^*=\Tr YY^*=1$. Let
$\alpha,\beta,\eta,\delta,K\ge0$ satisfy
\[
 \alpha+\beta+\eta+\delta=1,\qquad
 \alpha\delta-\beta\eta\ge K^2.
\]
For
\[
 T=\begin{pmatrix}
 \alpha XX^*+\beta YY^*&KXY^*\\
 KYX^*&\eta XX^*+\delta YY^*
 \end{pmatrix},
 \quad U_0=\alpha+\beta,\quad U_1=\eta+\delta,\quad
 u=U_0U_1-K^2,
\]
one has
\begin{equation}\label{eq:block}
 S(T)\ge(\alpha+\eta)S(XX^*)+(\beta+\delta)S(YY^*)+g(u).
\end{equation}
\end{theorem}
\begin{proof}
The assumptions ensure positivity of $T$: its coherent part with diagonal blocks $\alpha XX^*,\delta YY^*$ is positive because $\alpha\delta\ge K^2$, and the remaining diagonal blocks are positive.

Assume first $U_0U_1>0$ and $u>0$. With $U_0,U_1,K$ fixed, the constraint becomes
\[
 \beta U_1+\eta U_0\le u.
\]
In the coefficient order $(\alpha,\beta,\eta,\delta)$ the feasible triangle has vertices
\begin{align*}
 A&=(K^2/U_1,u/U_1,0,U_1),\\
 B&=(U_0,0,u/U_0,K^2/U_0),\\
 C&=(U_0,0,0,U_1).
\end{align*}
The original point has respective weights
\[
 w_A=\beta U_1/u,\qquad w_B=\eta U_0/u,\qquad
 w_C=1-w_A-w_B.
\]
They are nonnegative. Each vertex retains the same row masses and off-diagonal scalar $K$, hence the same $u$.

At $A$, write
\[
 T_A=M_AM_A^*,\qquad
 M_A=\begin{pmatrix}\sqrt{\alpha}X&\sqrt{\beta}Y\\
                     \sqrt{\delta}Y&0\end{pmatrix}.
\]
Here $\sqrt{\alpha\delta}=K$ and $\beta\delta=u$, so Lemma~\ref{lem:triangular} gives \eqref{eq:block} at this vertex. At $B$ use
\[
 M_B=\begin{pmatrix}\sqrt{U_0}X&0\\
 \sqrt{\delta}Y&\sqrt{\eta}X\end{pmatrix}.
\]
It satisfies $M_BM_B^*=T_B$; interchange its block rows to apply the triangular lemma, with corner-mass product $U_0\eta=u$.

At $C$, apply Lemma~\ref{lem:controlled} to
$\rho=\bigl(\begin{smallmatrix}U_0&K\\K&U_1\end{smallmatrix}\bigr)$ and the pair $X,Y$. The scalar state has determinant $u$ and entropy $g(u)$, giving the desired bound. Finally, entropy concavity applied to
$T=w_AT_A+w_BT_B+w_CT_C$ preserves $g(u)$ and the affine weights of the two block entropies.

If a row mass vanishes, $K=0$ and the result is ordinary entropy concavity in the remaining block. If $u=0$ and both row masses are positive, the constraint forces $\beta=\eta=0$; Lemma~\ref{lem:controlled} applies with a pure $\rho$. This also establishes all boundaries.
\end{proof}
The estimate is sharp for every allowed scalar coefficient tuple: setting $X=Y$ gives $T=\rho\otimes XX^*$, so equality holds. The equality boundary $\alpha\delta-\beta\eta=K^2$ is exactly the two-vertex thermal argument in \cite[Theorem 3.1]{openai}; the third vertex is needed to cover the inequality region.

\section{Population-resolved bound and capacity}
\begin{lemma}\label{lem:step}
Let $\mathcal N$ be any finite-dimensional channel and
$|\psi\rangle=\sqrt{1-p}|0\rangle|\psi_0\rangle+
\sqrt p|1\rangle|\psi_1\rangle$, where the nonzero branches are normalized. Then
\begin{equation}\label{eq:step}
 S((\Phi\otimes\mathcal N)(\psi))
 \ge e(p)+(1-p)S(\mathcal N(\psi_0))+pS(\mathcal N(\psi_1)).
\end{equation}
No orthogonality of the conditional vectors is required.
\end{lemma}
\begin{proof}
For Kraus operators $L_j$ of $\mathcal N$, form the matrices with columns
$X_j=L_j\psi_0$ and $Y_j=L_j\psi_1$. Their squared Hilbert--Schmidt norms are one. The output has the form in Theorem~\ref{thm:block}, with
\[
 \alpha=(1-p)(1-a),\quad\beta=pb,\quad
 \eta=(1-p)a,\quad\delta=p(1-b),\quad
 K=\lambda\sqrt{p(1-p)}.
\]
Here $\alpha\delta-\beta\eta=kp(1-p)\ge K^2$,
$\alpha+\eta=1-p$, and $u=v(p)$. Apply \eqref{eq:block}. At $p=0,1$ only one product branch remains and the formula follows directly.
\end{proof}

For clarity, the convex envelope can be written as
\begin{equation}\label{eq:roof}
 f(p)=\min_{\substack{0\le q,r,s\le1\\(1-s)q+sr=p}}
       \{(1-s)e(q)+se(r)\}.
\end{equation}
Indeed the minimum of $\int e\,d\mu$ over probability measures with mean $p$ is attained by compactness. This minimum is the lower convex envelope. An extreme minimizing measure has at most two support points: weights on three distinct support points admit a nonzero signed perturbation preserving total mass and mean. Equivalently, each lower supporting line of the convex hull of the graph needs at most two contact points. Continuity of $e$ includes endpoint and one-point cases.

\begin{proof}[Proof of Theorem~\ref{thm:main}]
Refining any signal ensemble into pure states does not decrease its Holevo information. A pure qubit with population $p$ has output entropy $e(p)$, independently of its phase. If an ensemble has mean population $\bar p$, dephasing its average output increases entropy, and convexity of $f\le e$ gives
\[
 (\ln2)\chi(\text{ensemble})
 \le h(a+k\bar p)-\sum_\ell w_\ell e(p_\ell)
 \le h(a+k\bar p)-f(\bar p).
\]
Conversely choose $q,r,s$ attaining \eqref{eq:roof}. At population $q$ use
$\sqrt{1-q}|0\rangle\pm\sqrt q|1\rangle$, each with probability $(1-s)/2$; at population $r$ use the analogous pair, each with probability $s/2$. The average is diagonal with population $p$, and the average output entropy is $f(p)$. Maximizing over $p$ proves the single-use formula and the four-signal bound.

For a pure ensemble $\{w_\ell,\psi_\ell\}$ on the input of $\Phi\otimes\mathcal N$, apply Lemma~\ref{lem:step} to every signal. The branch outputs of $\mathcal N$ with probabilities
$w_\ell(1-p_\ell),w_\ell p_\ell$ form a valid ensemble. Its average is exactly the second output marginal $\bar\omega_B$. Subadditivity bounds the entropy of the full average output by
$h(a+k\bar p)+S(\bar\omega_B)$. Subtracting the average of \eqref{eq:step} yields
\[
 (\ln2)\chi(\text{joint ensemble})
 \le h(a+k\bar p)-f(\bar p)+(\ln2)\chi(\mathcal N).
\]
The single-use formula and product ensembles prove \eqref{eq:chiadd}. The same one-step inequality bounds every pure output entropy below by
$\min_p e(p)+\Smin(\mathcal N)$; entropy concavity handles mixed inputs, and products attain the bound. This proves \eqref{eq:sadd}.

Iteration of \eqref{eq:chiadd} gives
$\chi(\Phi^{\otimes n}\otimes\mathcal N^{\otimes n})
=n\chi(\Phi)+\chi(\mathcal N^{\otimes n})$.
Taking the regularized capacity proves \eqref{eq:cadd}, including $C(\Phi)=\chi(\Phi)$.
\end{proof}


\section{Entanglement formation in the dilation range}\label{sec:formation}
For any finite channel $\mathcal M$ define its convex roof of output entropy,
in nats, by
\[
 \widehat S_{\mathcal M}(\rho)=
 \min_{\rho=\sum_jw_j\psi_j}\sum_jw_j S(\mathcal M(\psi_j)),
\]
where the $\psi_j$ are pure states. If $V:\mathbb C^2\to B\otimes E$
is a Stinespring isometry of $\Phi$, then
\begin{equation}\label{eq:msw}
 F_\Phi(\sigma):=\widehat S_\Phi(\sigma)
       =E_F(V\sigma V^*) .
\end{equation}
Here $E_F$ denotes bipartite state entanglement of formation in nats.
Equation~\eqref{eq:msw} is the standard channel/state correspondence
\cite{msw}. Each pure component of a state supported on $\operatorname{ran}V$
lies in that range and therefore pulls back to a pure qubit; this also gives
a direct proof of the equality.

An elementary observation strengthens the population argument. Every
qubit state $\sigma=\bigl(\begin{smallmatrix}1-p&z\\\bar z&p\end{smallmatrix}\bigr)$
is a mixture of pure states with the same population $p$: two opposite
phases, with suitably unequal weights, realize every $|z|\le\sqrt{p(1-p)}$.
Their output entropies are all $e(p)$. Thus
\begin{equation}\label{eq:roofupper}
 f(p)\le F_\Phi(\sigma)\le e(p).
\end{equation}
The upper bound, rather than equality with $f$, is what is needed next.

\begin{theorem}[Strong convex-roof superadditivity]\label{thm:roofadd}
For every finite-dimensional channel $\mathcal N$ and every joint input state
$\rho_{AA'}$,
\begin{equation}\label{eq:roofadd}
 \widehat S_{\Phi\otimes\mathcal N}(\rho)
 \ge \widehat S_\Phi(\rho_A)+\widehat S_{\mathcal N}(\rho_{A'}).
\end{equation}
The same assertion holds with $\Phi$ replaced by any complementary channel
$\Phi^c$.
\end{theorem}
\begin{proof}
For a pure $\rho=\psi$, Lemma~\ref{lem:step} bounds its output entropy by
$e(p)$ plus the average output entropies of the conditional $A'$ branches.
By \eqref{eq:roofupper}, the first term is at least
$\widehat S_\Phi(\psi_A)$. Those branches decompose $\psi_{A'}$, so their
average is at least $\widehat S_{\mathcal N}(\psi_{A'})$.
For mixed $\rho$, apply this to a minimizing pure decomposition and use
convexity of both marginal roofs.

For the complement, let $\mathcal N^c$ be a complementary channel to
$\mathcal N$. For pure $\psi$, purity of the complete Stinespring output gives
\[
 S((\Phi^c\otimes\mathcal N)(\psi))
   =S((\Phi\otimes\mathcal N^c)(\psi)).
\]
Apply Lemma~\ref{lem:step} to the right side. Its conditional branches are
pure, so $S(\mathcal N^c(\psi_j))=S(\mathcal N(\psi_j))$.
Thus the same one-step inequality holds for $\Phi^c$. Its pure-output entropy
is again $e(p)$, and the preceding argument applies. This use of complement
duality is standard \cite[Theorem 1, equations (6) and (9)]{holevocomp};
see also \cite{conjugate}.
\end{proof}

For a fixed input $\sigma$, the maximal ensemble Holevo information is
$(S(\mathcal M(\sigma))-\widehat S_{\mathcal M}(\sigma))/\ln2$.
Subadditivity of the mean output entropy and \eqref{eq:roofadd} therefore
give its subadditivity at arbitrary joint input marginals, with equality
at product input states. This is stronger than the unconstrained maximum in
Theorem~\ref{thm:main}; it uses the established convex-roof framework of
\cite{msw}, not an equivalence asserted without its needed hypotheses.

\begin{corollary}[State entanglement cost]\label{cor:cost}
Let $\omega=V\sigma V^*$. For every finite-dimensional bipartite state
$\Omega_{BB'EE'}$ with $\Omega_{BE}=\omega$,
\[
 E_F(\Omega_{BB':EE'})\ge E_F(\omega_{B:E})+
                              E_F(\Omega_{B':E'}).
\]
Consequently, for every finite bipartite state $\tau$,
\begin{align*}
 E_F(\omega\otimes\tau)&=E_F(\omega)+E_F(\tau),\\
 E_C(\omega)&=F_\Phi(\sigma)/\ln2,\\
 E_C(\omega\otimes\tau)&=F_\Phi(\sigma)/\ln2+E_C(\tau).
\end{align*}
Here $E_C$ is asymptotic LOCC \emph{state} entanglement cost in ebits.
\end{corollary}
\begin{proof}
Positivity and the support of the first marginal imply that $\Omega$ is
supported on $\operatorname{ran}V\otimes B'\otimes E'$. Lift it through $V$
and use Theorem~\ref{thm:roofadd} with $\mathcal N=\operatorname{Tr}_{E'}$.
The three roofs are exactly the three formation quantities displayed.
For product states, product pure decompositions attain the reverse inequality.
Iteration gives $E_F(\omega^{\otimes n})=nE_F(\omega)$, and more generally
$E_F(\omega^{\otimes n}\otimes\tau^{\otimes n})=
nE_F(\omega)+E_F(\tau^{\otimes n})$.
The state-cost theorem $E_C(\rho)=\lim_nE_F(\rho^{\otimes n})/(n\ln2)$
\cite{hht} proves both cost statements.
\end{proof}
These are states in a Stinespring range, not generally Choi states. The result
does not evaluate entanglement cost for simulating the channel. At the GAD
boundary the same corollary follows from OpenAI's one-step inequality and
the standard correspondence; our block extension supplies the remaining
Davies region.

For completeness the required one-use formation functional has a scalar
moment description. This is a convenient version of the known axial
convex-roof problem \cite{hu}, not a claim that its geometry is new.
\begin{proposition}\label{prop:moment}
With $r(q)=\sqrt{q(1-q)}$ and
$\sigma=\bigl(\begin{smallmatrix}1-p&z\\\bar z&p\end{smallmatrix}\bigr)$,
\begin{equation}\label{eq:moment}
 F_\Phi(\sigma)=\min_\mu\left\{\int e(q)\,d\mu(q):
       \int q\,d\mu=p,\quad\int r(q)\,d\mu\ge|z|\right\},
\end{equation}
where $\mu$ ranges over probability measures on $[0,1]$.
In particular $F_\Phi(\operatorname{diag}(1-p,p))=f(p)$.
\end{proposition}
\begin{proof}
A pure decomposition of $\sigma$ induces a population measure; the triangle
inequality for its coherences gives the moment constraint. Conversely, given
such a measure, put $R=\int r\,d\mu$. At every population choose the same
average phase factor $z/R$, with the sign convention matching the upper-right
matrix entry. Two opposite phases with probabilities
$(1+|z|/R)/2,(1-|z|/R)/2$ realize that factor and leave the entropy integral
unchanged. If $R=0$, necessarily $z=0$ and the endpoint mixture suffices.
Compactness gives a minimizing measure. An extreme minimizer can have at
most three populations: four supported pieces admit a nonzero signed weight
perturbation preserving normalization, the population moment and the radius
moment. This also supplies a finite ensemble. For $z=0$ the constraint is
vacuous, reducing the program to the convex envelope.
\end{proof}
Formula~\eqref{eq:moment} is an optimization formula, not an elementary closed
form. For example, the sufficient condition $D\ge0$, $4AD\ge B^2$, with
$v(q)=A+Bq+Dq^2$, makes $e$ convex and then $F_\Phi(\sigma)=e(p)$ for every
coherence. This is the flat-roof regime already identified by Hellmund and
Uhlmann \cite[Theorem 2 and Sections IV.C--D]{hu}. One quick verification is
that $\sqrt v$ is a norm of an affine two-vector under this condition, while
$g(x^2)$ is convex and nondecreasing \cite{wootters,hu}. We do not claim this
criterion or the flat roof as new. Outside that regime, \eqref{eq:moment}
retains the coherence constraint even though additivity still holds.

\section{Classical capacity of the complementary channel}\label{sec:complement}
The complementary output is the environment output of a pure Stinespring
dilation, including purification of a mixed reservoir when necessary. It is
not automatically the experimentally accessible part of a thermal bath.
Its minimal dimension is at most four. To compute its entropy at a diagonal
input, choose unit vectors $u,v$ in an environment sector of dimension two,
and orthogonal vectors $e_+,e_-$ orthogonal to that sector, with
$\langle v,u\rangle=\lambda/\sqrt{(1-a)(1-b)}$. A dilation is
\begin{align*}
 V|0\rangle&=\sqrt{1-a}|0\rangle|u\rangle+
                         \sqrt a|1\rangle|e_+\rangle,\\
 V|1\rangle&=\sqrt b|0\rangle|e_-\rangle+
                         \sqrt{1-b}|1\rangle|v\rangle.
\end{align*}
The positivity constraint makes the prescribed inner product feasible.
If its denominator vanishes then $\lambda=0$ and the irrelevant overlap may
be chosen arbitrarily. Tracing the environment gives \eqref{eq:channel}.
For input $\operatorname{diag}(1-p,p)$ the environment spectrum is
\begin{equation}\label{eq:envspectrum}
 a(1-p),\quad bp,\quad
 \frac{w(p)+\sqrt{w(p)^2-4d(p)}}2,\quad
 \frac{w(p)-\sqrt{w(p)^2-4d(p)}}2,
\end{equation}
where
\[
 w(p)=(1-a)(1-p)+(1-b)p,\qquad
 d(p)=p(1-p)((1-a)(1-b)-\lambda^2).
\]
Indeed the remaining block is the mixture of $|u\rangle\langle u|$ and
$|v\rangle\langle v|$ with the indicated trace and determinant. Let $H_E(p)$
be the entropy in nats of the four probabilities in \eqref{eq:envspectrum}.

\begin{theorem}\label{thm:complement}
For every qubit Davies channel in \eqref{eq:region},
\[
 C(\Phi^c)=\chi(\Phi^c)=
       \frac1{\ln2}\max_{0\le p\le1}\{H_E(p)-f(p)\}.
\]
At most two opposite-phase population pairs attain the one-use maximum.
For every finite-dimensional $\mathcal N$,
\begin{align*}
 \chi(\Phi^c\otimes\mathcal N)&=\chi(\Phi^c)+\chi(\mathcal N),\\
 C(\Phi^c\otimes\mathcal N)&=\chi(\Phi^c)+C(\mathcal N),\\
 \Smin(\Phi^c\otimes\mathcal N)&=\min_p e(p)+\Smin(\mathcal N).
\end{align*}
\end{theorem}
\begin{proof}
For $U_\theta=\operatorname{diag}(1,e^{i\theta})$, let $W_\theta$ act as the
identity on $u,v$ and with phases $e^{-i\theta},e^{i\theta}$ on $e_+,e_-$.
Then $VU_\theta=(U_\theta\otimes W_\theta)V$. Hence $\Phi^c$ is phase
covariant, and averaging the input phase cannot decrease its output entropy.
At mean population $p$ that entropy is therefore at most $H_E(p)$, while
the mean pure-output entropy is at least $f(p)$. The same opposite-phase
ensemble used in Theorem~\ref{thm:main} attains both bounds.

The proof of Theorem~\ref{thm:roofadd} established the full one-step inequality
for $\Phi^c$. Repeating the branch-ensemble proof of Theorem~\ref{thm:main},
with $H_E$ replacing $h\circ t$, gives the arbitrary-partner Holevo and
minimum-entropy identities. Regularization gives the capacity identity.
\end{proof}
The function $H_E$ is concave, being entropy along an affine family of states;
therefore $H_E-f$ is concave. A sampled lower hull $f_N$ gives the same
one-sided discretization error as Proposition~\ref{prop:grid}, since only
the error in $f$ enters that argument. The limits provide useful checks:
an identity channel has a constant complement of capacity zero, whereas a
constant qubit-output channel has complementary capacity one bit.


\section{Parallel simulation of the complementary channel}\label{sec:simulation}
This section concerns the resources needed to simulate $\mathcal N=\Phi^c$,
with its full Stinespring output. Let
\[
 G_{\mathcal M}(\sigma)=E_F((\operatorname{id}_R\otimes\mathcal M)(\psi_\sigma)),
 \qquad E_F^{\rm ch}(\mathcal M)=\max_\sigma G_{\mathcal M}(\sigma),
\]
where $\psi_\sigma$ is a purification and $E_F$ uses nats. Reference isometries
do not affect this definition. The parallel simulation cost $E_C^{\rm par}$
is the asymptotic number of shared ebits per use required by LOCC protocols
whose simulated tensor-power channel converges in diamond norm; classical
communication is free. The theorem of Berta et al.\ gives
\begin{equation}\label{eq:channelcostreg}
 E_C^{\rm par}(\mathcal M)=\lim_{n\to\infty}
       \frac{E_F^{\rm ch}(\mathcal M^{\otimes n})}{n\ln2}.
\end{equation}
This operational theorem is external \cite[Theorem 12]{bcbw}. Sequential or
adaptive tests constitute a different task \cite{wilde}.

We recall two standard formation facts. The Kraus characterization is
\[
 G_{\mathcal M}(\sigma)=\inf_{\{L_j\}}
 \sum_j\operatorname{Tr}(L_j\sigma L_j^*)
 S\left(\frac{L_j\sigma L_j^*}{\operatorname{Tr}(L_j\sigma L_j^*)}\right),
\]
with zero branches omitted. Pure-state ensemble decompositions of the
purified output correspond to Kraus representations; one first argues on
the support of $\sigma$ and extends the representation or takes continuous
limits at singular inputs. For every fixed representation the entropy
perspective is concave in $\sigma$. Its infimum is concave as well: bound
each component of a mixture by its own infimum before taking the infimum
over the common representation. Thus $G_{\mathcal M}$ is concave. It is
continuous by continuity of finite-dimensional $E_F$ and of the canonical
purification constructed from $\sqrt\sigma$, so a maximum exists. Covariance
and local-unitary invariance then imply that $G_{\Phi^c}$ has a diagonal
maximizer, by averaging the input phase.

The other fact is subadditivity of channel formation
\cite[Appendix C]{tang}:
\begin{equation}\label{eq:chsub}
 E_F^{\rm ch}(\mathcal M\otimes\mathcal L)
 \le E_F^{\rm ch}(\mathcal M)+E_F^{\rm ch}(\mathcal L).
\end{equation}
For completeness, choose Kraus decompositions optimal, or arbitrarily close
to optimal, for the two marginal input states of a joint purified input.
Their product yields an admissible pure decomposition of the joint output.
For its classical Kraus labels $I,J$ and quantum outputs $B,B'$, entropy
subadditivity and conditioning give
\[
 S(BB'|IJ)\le S(B|IJ)+S(B'|IJ)
             \le S(B|I)+S(B'|J).
\]
Completeness of the other instrument makes each final term the corresponding
local Kraus objective. Maximizing over inputs proves \eqref{eq:chsub}.

\begin{theorem}[Complementary simulation]\label{thm:simulation}
For every Davies channel $\Phi$ in \eqref{eq:region} and every finite channel
$\mathcal M$,
\begin{align*}
 E_F^{\rm ch}(\Phi^c\otimes\mathcal M)
   &=E_F^{\rm ch}(\Phi^c)+E_F^{\rm ch}(\mathcal M),\\
 E_C^{\rm par}(\Phi^c)&=E_F^{\rm ch}(\Phi^c)/\ln2,\\
 E_C^{\rm par}(\Phi^c\otimes\mathcal M)
   &=E_F^{\rm ch}(\Phi^c)/\ln2+E_C^{\rm par}(\mathcal M).
\end{align*}
An exact scalar formula is
\begin{equation}\label{eq:simformula}
 E_C^{\rm par}(\Phi^c)=\frac1{\ln2}\max_{0\le p\le1} f_p(t(p)),
\end{equation}
where $f_p$ is the lower convex envelope of the pure-output entropy for the
reversed Davies parameters
\begin{equation}\label{eq:reverse}
 a_p=\frac{pb}{1-t},\qquad b_p=\frac{(1-p)a}{t},\qquad
 k_p=\frac{kp(1-p)}{t(1-t)},\qquad
 \lambda_p^2=\frac{\lambda^2p(1-p)}{t(1-t)},\quad t=t(p).
\end{equation}
Endpoint and constant pure-output exceptions are specified in the proof.
\end{theorem}
\begin{proof}
For $0<p<1$ and $0<t<1$, apply the dilation of Section~\ref{sec:complement}
to half of $\sqrt{1-p}|00\rangle+\sqrt p|11\rangle$. Grouping the resulting
pure vector by its $B$ output gives an isometry $W_p$ from a qubit to $R E$:
\begin{align*}
 W_p|0\rangle&=\frac{\sqrt{(1-p)(1-a)}|0\rangle|u\rangle+
                           \sqrt{pb}|1\rangle|e_-\rangle}{\sqrt{1-t}},\\
 W_p|1\rangle&=\frac{\sqrt{(1-p)a}|0\rangle|e_+\rangle+
                           \sqrt{p(1-b)}|1\rangle|v\rangle}{\sqrt t}.
\end{align*}
Its columns are orthonormal. Tracing $E$ gives a qubit channel $\Psi_p$ with
parameters \eqref{eq:reverse}; in particular $1-a_p-b_p=k_p$ and
$0\le\lambda_p^2\le k_p$. This is the key closure property. The
reference--environment output state is exactly
\[
 \omega_p=(\operatorname{id}_R\otimes\Phi^c)(\psi_p)
            =W_p\operatorname{diag}(1-t,t)W_p^*.
\]
Corollary~\ref{cor:cost} therefore gives strong formation additivity for
$\omega_p$ with every finite bipartite partner, and
$G_{\Phi^c}(\operatorname{diag}(1-p,p))=f_p(t)$.

Choose a diagonal input maximizing $G_{\Phi^c}$ and an optimizing purified
input for $\mathcal M$, with output state $\tau$. Product inputs, together
with the established state additivity, give
\[
 E_F^{\rm ch}(\Phi^c\otimes\mathcal M)
 \ge E_F(\omega_p\otimes\tau)
 =E_F^{\rm ch}(\Phi^c)+E_F^{\rm ch}(\mathcal M).
\]
Equation~\eqref{eq:chsub} proves equality. Iteration and
\eqref{eq:channelcostreg} prove both operational identities and
\eqref{eq:simformula}.

At $p=0,1$ the reference is pure and $G=0$; a zero maximizing formation value
also gives the product lower bound by local discard. If $t=0$ or $1$ at an
interior $p$, the original channel is constant with pure output. Its
complement is isometric, $G(\operatorname{diag}(1-p,p))=h(p)$ and its cost
is one ebit. Pure maximal entanglement gives its strong product lower bound.
These cases define the missing values of \eqref{eq:simformula}; all other
boundary cases follow from the displayed construction or continuity.
\end{proof}
This proof uses an optimizing reference--output state in a different
Stinespring range; complement duality alone does not transfer simulation
cost. The construction is the standard normalized reverse-channel
(weighted-transpose or Petz) construction \cite{barnumknill}, with its
explicit Davies parameters verified here. At the
GAD boundary the same deduction uses the antecedent one-step inequality
\cite{openai}. The added dephasing and nonunital parameters are the extension.

\begin{corollary}[An elementary rate]\label{cor:closedcost}
Put $x=\lambda^2$ and
\[
 c(a,b)=k+2ab-2\sqrt{ab(k+ab)}.
\]
If $c(a,b)\le x\le k$, then
\[
 E_C^{\rm par}(\Phi^c)=h\left(\frac{1+\lambda}{2}\right)/\ln2.
\]
A Bell input maximizes channel formation in this region.
\end{corollary}
\begin{proof}
The threshold $c$ is the known Hellmund--Uhlmann flat-concurrence threshold,
a sufficient condition for the entropy roof used here
\cite[Equation (29) and Section V.B]{hu}. With $L=p(1-p)/(t(1-t))$, the reverse parameters
satisfy $k_p=Lk$, $a_pb_p=Lab$, $\lambda_p^2=Lx$, and hence
$c(a_p,b_p)=Lc(a,b)$. Their entropy roof is therefore flat in the stated
region. Direct substitution yields
\[
 f_p(t)=e_p(t)=g((1-x)p(1-p)).
\]
The maximum is at $p=1/2$, proving the formula. The exceptional constant
pure-output case has the same value one; continuity covers the remaining
endpoints.
\end{proof}
The threshold itself is an antecedent, not a new convex-roof criterion or
a necessary boundary for entropy flatness or for this simulation formula.
For example, $a=0.1,b=0.3,\lambda^2=0.5$ has complementary parallel cost
about $0.600876037$ ebits per use. This determines the complementary
simulation rate, not the simulation cost of $\Phi$.

There is also a sufficient low-coherence regime with a simpler optimization.
Define the classical conditional entropy
\[
 J(p)=h(p)+(1-p)h(a)+p h(b)-h(a+kp).
\]
The endpoint mixture in every reversed roof gives $f_p(t)\le J(p)$, by the
classical entropy chain rule. If $p_0$ maximizes $J$, is interior, and has
$0<t_0=t(p_0)<1$, then
\begin{equation}\label{eq:plateau}
 \lambda^2\le\frac{k^2p_0(1-p_0)}{t_0(1-t_0)}
 \quad\Longrightarrow\quad
 E_C^{\rm par}(\Phi^c)=J(p_0)/\ln2.
\end{equation}
Indeed, the quadratic coefficient of the reversed determinant at $p_0$ is
$L(\lambda^2-Lk^2)\le0$. Since $g$ is increasing and concave, $e_{p_0}$ is
concave. Its convex envelope is its endpoint chord, attaining the upper
bound $J(p_0)$. Equation~\eqref{eq:plateau} is sufficient; no complete
transition diagram is asserted.

\begin{proposition}[Strict separation from Choi-state cost]\label{prop:choigap}
For $a=1/10,b=3/10,\lambda^2=1/10$, let $\mathcal N=\Phi^c$ and let
$J_{\mathcal N}$ be its normalized Choi state. Then
\[
 E_C^{\rm par}(\mathcal N)>E_C(J_{\mathcal N}).
\]
In particular a Bell input does not maximize $E_F^{\rm ch}(\mathcal N)$.
\end{proposition}
\begin{proof}
At $p=1/2$ the reversed parameters have $a_p=1/4$, $b_p=1/8$,
$k_p=5/8$ and $\lambda_p^2=5/48$. The quadratic coefficient is
$5/48-25/64=-55/192<0$, so it remains negative in an open interval of $p$.
There the endpoint-chord argument gives $G_{\mathcal N}(p)=J(p)$ exactly.
But
\[
 J'(1/2)=h(3/10)-h(1/10)-\tfrac35\ln(3/2)
       =\tfrac1{10}\ln\left(\frac{2^6 3^9}{7^7}\right)>0,
\]
because $1259712>823543$. Thus an input with $p>1/2$ has strictly larger
formation than the Choi input. Theorem~\ref{thm:simulation} identifies the
channel cost with its maximum formation, whereas Corollary~\ref{cor:cost},
applied to $W_{1/2}$, identifies the Choi-state cost with its formation.
\end{proof}
Consequently this channel cannot be simulated exactly by LOCC using one
copy of its own Choi state per call: preparing many such states at their
state cost would contradict the strict inequality. Other resource states
are not excluded. The example has $\lambda^2>ab$ and is nonunital, so it
does not merely restate the covariant epolarizing example of \cite{wilde}.
Approximate sampled values are $0.704192652$ ebits at the Bell input and
$0.704712865$ near $p=0.516949108$; the strict proof does not use these values.

\section{Finite-block consequences}
The preceding result also applies to nonidentical channels in \eqref{eq:region}. Write $e_i,f_i,t_i$ for the corresponding functions of channel $\Phi_i$.

\begin{corollary}\label{cor:multi}
For every state $\rho$ on $n$ qubits, with marginal excitation probabilities $p_i$,
\[
 S((\Phi_1\otimes\cdots\otimes\Phi_n)(\rho))
 \ge \sum_{i=1}^n f_i(p_i).
\]
\end{corollary}
\begin{proof}
For a pure input split at the first site and use \eqref{eq:step}. Apply induction on each conditional remaining state and combine its populations using convexity of each $f_i$. At the first site $e_1\ge f_1$. A pure-state decomposition, entropy concavity and convexity of the $f_i$ give the mixed-input statement.
\end{proof}

This yields an exact accounting of the Holevo deficit. Let a pure ensemble have output states $\omega_\ell$, mean $\bar\omega$, populations $p_{i\ell}$ and means $\bar p_i$. Put
\[
 c_i=(\ln2)\chi(\Phi_i),\quad
 \sigma=\bigotimes_i\operatorname{diag}\bar\omega_i,\quad
 R_\ell=S(\omega_\ell)-\sum_i f_i(p_{i\ell})\ge0.
\]
Using quantum relative entropy in nats, direct subtraction gives
\begin{align}\label{eq:deficit}
 \sum_i c_i-(\ln2)\chi(\text{ensemble})
 ={}&\sum_i[c_i-h(t_i(\bar p_i))+f_i(\bar p_i)]\notag\\
 &+D(\bar\omega\Vert\sigma)+\sum_\ell w_\ell R_\ell\\
 &+\sum_i\left[\sum_\ell w_\ell f_i(p_{i\ell})-f_i(\bar p_i)\right].\notag
\end{align}
Every term is nonnegative. In particular,
$D(\bar\omega\Vert\sigma)=\sum_i h(t_i(\bar p_i))-S(\bar\omega)$.
The support convention is harmless: a zero marginal basis probability gives zero support of $\bar\omega$ on that sector. Equation \eqref{eq:deficit} is a consequence of the population bound and standard relative-entropy identities. It extends the usual finite-block deficit method, including that used for amplitude damping in \cite{tang}; it does not imply uniqueness of optimizers or a uniform positive gap.

\section{Scalar computation and examples}
The envelope in \eqref{eq:roof} can be evaluated by the lower convex hull of the sampled curve. The next bound makes the discretization error explicit, separately from floating-point error.

\begin{proposition}\label{prop:grid}
Suppose $e$ is $L$-Lipschitz on $[0,1]$. For a uniform grid of mesh $1/N$, let $f_N$ be the lower convex envelope of its sampled points, linearly extended between hull vertices. If
\[
 C_N=\frac1{\ln2}\max_p\{h(a+kp)-f_N(p)\},
\]
then
\[
 C_N\le C(\Phi)\le C_N+\frac{L}{2N\ln2}.
\]
\end{proposition}
\begin{proof}
The piecewise-linear interpolant $\widetilde e_N$ satisfies
$|\widetilde e_N-e|\le L/(2N)$: at relative position $\theta$ in a cell the error is at most $2L\theta(1-\theta)/N$.
Taking lower convex envelopes preserves uniform error bounds. Moreover $f_N\ge f$, because mixtures of sampled populations form a subset of all mixtures; and $\co\widetilde e_N=f_N$. Thus
$0\le f_N-f\le L/(2N)$, which gives the claim.
\end{proof}

For parameters with $m=\min_p v(p)>0$, write
\[
 v(p)=A+Bp+Dp^2,\quad
 A=a(1-a),\quad B=k(1-2a)-\lambda^2,\quad D=\lambda^2-k^2.
\]
The minimum is found from the endpoints and any interior vertex of this quadratic. A valid Lipschitz constant is
\[
 L=\max\{|B|,|B+2D|\}\,g'(m),\qquad
 g'(m)=\frac{2\operatorname{atanh}\sqrt{1-4m}}{\sqrt{1-4m}},
\]
with limit $2$ at $m=1/4$. The function $g'$ is nonincreasing; this follows by differentiating $\operatorname{atanh}x/x$. When $m=0$ one can instead use uniform continuity; the displayed finite-$L$ bound is not asserted at that boundary.

For $k>0$, on a hull segment of slope $s$, the maximizing population is the projection onto that segment of
\[
 p=\frac{(1+\exp(s/k))^{-1}-a}{k}.
\]
This follows by setting $k\ln((1-t)/t)=s$ and using concavity. The case $k=0$ has zero capacity.

A strict interior example is
\[
 a=\frac1{10},\quad b=\frac3{10},\quad k=\frac35,\quad
 \lambda^2=\frac12.
\]
It is nonunital, has additional dephasing, and is not entanglement breaking. In fact the Choi partial transpose is positive exactly when $\lambda^2\le ab$, which fails here. Its Choi rank is four, and $\min v=5/56>0$, so it admits no pure output. The example is therefore not covered merely by the unital, entanglement-breaking, pure-output or $\lambda^2=k$ boundary cases.

\begin{center}
\begin{tabular}{rrrrr}
\toprule
$a$&$b$&$\lambda^2$&capacity (bits, approximate)&mean population\\
\midrule
0.10&0.30&0.60&0.503934607&0.468332\\
0.10&0.30&0.50&0.407849860&0.470350\\
0.10&0.30&0.10&0.296671803&0.471876\\
0.02&0.30&0.514624&0.436551160&0.441065\\
\bottomrule
\end{tabular}
\end{center}
These are reproducible floating-point illustrations on a grid with $N=200000$, not machine-certified intervals. Proposition~\ref{prop:grid}'s exact-arithmetic error bounds are below $3.8\times10^{-6}$ bits for these cases, before rounding error. The last example's sampled hull optimizer mixes population $0$ with approximately $0.958865$, illustrating why a single population pair need not suffice. No exact three-signal necessity or global transition boundary is inferred from that numerical observation. For the same four channels, a separate computation with $N=60000$ gives complementary capacities approximately $0.492274059$, $0.562095952$, $0.906804474$ and $0.469017072$ bits. Direct Stinespring matrices reproduce \eqref{eq:envspectrum} and the phase covariance in 500 random cases, with maximum entry or eigenvalue error below $5.0\times10^{-16}$. These are floating-point illustrations, not machine-certified intervals. For $a=0.02,b=0.30,\lambda^2=0.514624$, a sampled version of \eqref{eq:moment} at $p=0.44$ gives state formation costs about $0.466960$ ebits at $z=0$, $0.467338$ at $z=0.25$, and $0.469937$ at $z=0.45$. The coherence constraint matters; the scalar roof $f(p)$ alone does not describe all states. The linear program uses 20001 population samples and reconstructs a feasible phase ensemble up to numerical tolerance; no continuous optimization certificate is claimed.

\section{Verification, dependence and scope}
The proof uses Lemma~\ref{lem:triangular} as an external result and gives all subsequent matrix and channel arguments explicitly. Some minimum-output-entropy subregions are also covered by Fukuda's composition criterion \cite{fukuda}; that criterion is not a generic closure theorem for unrestricted Holevo additivity under composition. A root-agent exploratory screen tested 16000 random rectangular block instances without a violation. Separately, an Astra candidate audit tested 9000 block instances, their 27000 vertex instances, 1200 channel outputs with arbitrary finite partner samples, and 300 ensembles. Its smallest residual was of order $10^{-14}$ in equality cases. These floating-point tests are error-finding aids, not proofs.

A tempting alternative argument---that extra dephasing must increase the block entropy by at least the scalar entropy increase---failed a numerical check and is not used. The valid argument preserves the scalar entropy across the three vertices and invokes concavity only after their separate estimates.

The present family excludes the permitted phase-covariant region $\lambda^2>1-a-b$ and channels with negative longitudinal multiplier. We determine the state costs in Corollary~\ref{cor:cost} and the complementary parallel simulation cost in Theorem~\ref{thm:simulation}. We do not determine the original channel's simulation cost, adaptive simulation cost, quantum or private capacity, interacting multiqubit Davies dynamics, or finite-error coding rates. Memoryless use is essential to the operational interpretation. No claim of exhaustive literature coverage or worldwide priority accompanies the stated extension; its mathematical scope is explicit and its immediate antecedents are identified.

\paragraph{AI assistance and review.}
The author directed an AI-assisted research process. Codex developed the extension, proofs, exposition and scalar calculations. A separately dispatched reviewer requested as GPT-6 Astra with xhigh effort audited frozen proof dossiers and an earlier capacity-only manuscript. The reviewer has prior involvement and is not blind to the work. A subsequent integrated version was also reviewed. The present version further adds complementary parallel simulation, reverse-channel closure and the strict Choi-cost separation; every manuscript assessment must identify its exact input version, and an earlier assessment does not automatically apply to these additions. Model runtime identity was not independently attested. No human peer review or predetermined numerical grade is claimed.

\paragraph{Availability.}
The supplement contains the proof dossier, numerical scripts and reports, and source-retrieval metadata. Original implementation code is released under CC0; the manuscript is offered under CC BY 4.0. External papers retain their authorship and licences. The scalar scripts do not implement directed-rounding interval certification.

\begin{thebibliography}{99}
\small\setlength{\itemsep}{3pt}
\bibitem{roga} W. Roga, M. Fannes and K. \.{Z}yczkowski, \emph{Davies maps for qubits and qutrits}, Reports on Mathematical Physics 66, 311--329 (2010). \href{https://arxiv.org/abs/0911.5607}{arXiv:0911.5607v2}, Sections 3.1--3.3.
\bibitem{tang} Z. Tang, C. Zhu, G. Bai and X. Wang, \emph{Classical capacity and entanglement cost of the amplitude damping channel}, \href{https://arxiv.org/abs/2609.28592}{arXiv:2609.28592v1} (2026), Theorem 4.2 and Section 5.
\bibitem{openai} OpenAI, \emph{Classical capacity and entropy inequalities for generalized amplitude damping}, 24 September 2026, Theorems 2.1 and 3.1, Corollary 1.2. \href{https://github.com/openai/math/tree/main/preprints/Classical-capacity-and-entropy-inequalities-for-generalized-amplitude-damping-September-24-2026}{OpenAI mathematics repository}. Retrieved 17 October 2026; PDF SHA256 recorded in the supplement.
\bibitem{ssa} E. H. Lieb and M. B. Ruskai, \emph{Proof of the strong subadditivity of quantum-mechanical entropy}, Journal of Mathematical Physics 14, 1938--1941 (1973). \href{https://doi.org/10.1063/1.1666274}{doi:10.1063/1.1666274}.
\bibitem{king} C. King, \emph{Additivity for unital qubit channels}, Journal of Mathematical Physics 43, 4641--4653 (2002). \href{https://doi.org/10.1063/1.1500791}{doi:10.1063/1.1500791}.
\bibitem{shor} P. W. Shor, \emph{Additivity of the classical capacity of entanglement-breaking quantum channels}, Journal of Mathematical Physics 43, 4334--4340 (2002). \href{https://arxiv.org/abs/quant-ph/0201149}{arXiv:quant-ph/0201149}.
\bibitem{leditzky} F. Leditzky, E. Kaur, N. Datta and M. M. Wilde, \emph{Approaches for approximate additivity of the Holevo information of quantum channels}, Physical Review A 97, 012332 (2018). \href{https://arxiv.org/abs/1709.01111}{arXiv:1709.01111}.
\bibitem{holevo} A. S. Holevo, \emph{The capacity of the quantum channel with general signal states}, IEEE Transactions on Information Theory 44, 269--273 (1998). \href{https://arxiv.org/abs/quant-ph/9611023}{arXiv:quant-ph/9611023}.
\bibitem{sw} B. Schumacher and M. D. Westmoreland, \emph{Sending classical information via noisy quantum channels}, Physical Review A 56, 131--138 (1997). \href{https://doi.org/10.1103/PhysRevA.56.131}{doi:10.1103/PhysRevA.56.131}.
\bibitem{amosov} G. G. Amosov, \emph{On estimating the output entropy of a tensor product of the quantum phase-damping channel with an arbitrary channel}, Problems of Information Transmission 49, 224--231 (2013). \href{https://arxiv.org/abs/1301.2886}{arXiv:1301.2886}, Proposition 5.
\bibitem{rogathesis} W. Roga, \emph{Entropy of quantum channel in the theory of quantum information}, PhD thesis, \href{https://arxiv.org/abs/1108.5065v3}{arXiv:1108.5065v3} (2011), Sections 5.6 and 6.
\bibitem{msw} K. Matsumoto, T. Shimono and A. Winter, \emph{Remarks on additivity of the Holevo channel capacity and of the entanglement of formation}, Communications in Mathematical Physics 246, 427--442 (2004). \href{https://arxiv.org/abs/quant-ph/0206148}{arXiv:quant-ph/0206148v4}, Sections 4--5.
\bibitem{hht} P. M. Hayden, M. Horodecki and B. M. Terhal, \emph{The asymptotic entanglement cost of preparing a quantum state}, Journal of Physics A 34, 6891--6898 (2001). \href{https://arxiv.org/abs/quant-ph/0008134}{arXiv:quant-ph/0008134}, Theorem 1.
\bibitem{holevocomp} A. S. Holevo, \emph{On complementary channels and the additivity problem}, \href{https://arxiv.org/abs/quant-ph/0509101v3}{arXiv:quant-ph/0509101v3} (2005), Theorem 1.
\bibitem{conjugate} C. King, K. Matsumoto, M. Nathanson and M. B. Ruskai, \emph{Properties of conjugate channels with applications to additivity and multiplicativity}, Markov Processes and Related Fields 13, 391--423 (2007). \href{https://arxiv.org/abs/quant-ph/0509126v3}{arXiv:quant-ph/0509126v3}.
\bibitem{hu} M. Hellmund and A. Uhlmann, \emph{Concurrence and entanglement entropy of stochastic 1-qubit maps}, Physical Review A 79, 052319 (2009). \href{https://arxiv.org/abs/0903.1340}{arXiv:0903.1340}, Sections IV.C--D and V.
\bibitem{wootters} W. K. Wootters, \emph{Entanglement of formation of an arbitrary state of two qubits}, Physical Review Letters 80, 2245--2248 (1998). \href{https://arxiv.org/abs/quant-ph/9709029}{arXiv:quant-ph/9709029}.
\bibitem{fukuda} M. Fukuda, \emph{Extending additivity from symmetric to asymmetric channels}, Journal of Physics A 38, L753--L758 (2005). \href{https://arxiv.org/abs/quant-ph/0505022}{arXiv:quant-ph/0505022v4}.
\bibitem{bcbw} M. Berta, F. G. S. L. Brand\~ao, M. Christandl and S. Wehner, \emph{Entanglement cost of quantum channels}, IEEE Transactions on Information Theory 59, 6779--6795 (2013). \href{https://arxiv.org/abs/1108.5357v3}{arXiv:1108.5357v3}, Theorem 12.
\bibitem{wilde} M. M. Wilde, \emph{Entanglement cost and quantum channel simulation}, Physical Review A 98, 042338 (2018). \href{https://arxiv.org/abs/1807.11939v3}{arXiv:1807.11939v3}, Sections II--IV.
\bibitem{barnumknill} H. Barnum and E. Knill, \emph{Reversing quantum dynamics with near-optimal quantum and classical fidelity}, Journal of Mathematical Physics 43, 2097--2106 (2002). \href{https://arxiv.org/abs/quant-ph/0004088v1}{arXiv:quant-ph/0004088v1}, Equation (4), for the normalized reversal construction.
\end{thebibliography}
\end{document}
